\section*{अध्याय \ref{ch:b3:complex-mechanisms} --- संकुलों की अभिक्रिया क्रियाविधियाँ}

\begin{solution}{exo:b3:complex-mechanisms:1}\emergencystretch=3em
अक्रिय: $\ce{[Cr(H2O)6]^{3+}}$ ($d^3$) और $\ce{[Co(NH3)6]^{3+}}$ (निम्न-चक्रण $d^6$), बड़ी लिगैंड-क्षेत्र
सक्रियण ऊर्जाएँ। \omterm{def:b3:complex-mechanisms:labile}{विनिमयशील}: $\ce{[Cr(H2O)6]^{2+}}$ (उच्च-चक्रण $d^4$) और $\ce{[Cu(H2O)6]^{2+}}$ ($d^9$),
लंबे, दुर्बल अक्षीय आबंधों के साथ यान--टेलर विकृत; $\ce{[Zn(H2O)6]^{2+}}$ ($d^{10}$), कोई लिगैंड-क्षेत्र
अवरोध नहीं।
\end{solution}

\begin{solution}{exo:b3:complex-mechanisms:2}
जब $k_2[\mathrm Y] \gg k_{-1}[\mathrm X]$, तो हर $k_2[\mathrm Y]$ है और $v = k_1[\mathrm{ML_5X}]$; जब $k_{-1}[\mathrm X] \gg k_2[\mathrm Y]$, तो यह
$k_{-1}[\mathrm X]$ है और दूसरा रूप प्राप्त होता है (निर्गामी लिगैंड द्वारा संदमित
एक पूर्व-साम्य)।
\end{solution}

\begin{solution}{exo:b3:complex-mechanisms:3}
D: $\Delta V^{\ddagger} > 0$ और $\Delta^{\ddagger}S^\circ > 0$ (एक लिगैंड निकलता है)। A: दोनों ऋणात्मक (संक्रमण अवस्था में
एक लिगैंड बँधा होता है)।
\end{solution}

\begin{solution}{exo:b3:complex-mechanisms:4}
$cis$-$\ce{[PtCl2(NH3)2]}$, $\ce{[PtCl4]^{2-}}$ से; $trans$-$\ce{[PtCl2(NH3)2]}$, $\ce{[Pt(NH3)4]^{2+}}$ से (\ce{Cl-} का
\omterm{def:b3:complex-mechanisms:trans-effect}{ट्रांस प्रभाव} \ce{NH3} से अधिक है)।
\end{solution}

\begin{solution}{exo:b3:complex-mechanisms:5}
$k_{\mathrm{obs}} = 2.0 \times \num{3.0e4} \times 0.10/1.2 = \qty{5.0e3}{s^{-1}}$; \qty{1.0}{mol/L} पर $\num{6.0e4}/3.0 = \qty{2.0e4}{s^{-1}}$; सीमा
$k_i = \qty{3.0e4}{s^{-1}}$।
\end{solution}

\begin{solution}{exo:b3:complex-mechanisms:6}
$\Delta V^{\ddagger} = -RT\ln2/\Delta p = -8.314 \times 298 \times 0.693/\num{99.9e6} = \qty{-1.7e-5}{m^3.mol^{-1}} = \qty{-17}{cm^3.mol^{-1}}$:
सहयोजी।
\end{solution}

\begin{solution}{exo:b3:complex-mechanisms:7}
फ़ॉस्फ़ीन, एक प्रबल $\sigma$ दाता और $\pi$ ग्राही, अपने ट्रांस स्थित आबंध को दुर्बल करता है (\omterm{def:b3:complex-mechanisms:trans-effect}{ट्रांस
अनुप्रभाव}, लंबा आबंध) और उसे \omterm{def:b3:complex-mechanisms:labile}{विनिमयशील} बनाता है (\omterm{def:b3:complex-mechanisms:trans-effect}{ट्रांस प्रभाव}): वही क्लोराइड
पहले प्रतिस्थापित होता है।
\end{solution}

\begin{solution}{exo:b3:complex-mechanisms:8}
सेतु लिगैंड अंततः ऑक्सीकृत धातु पर स्थानांतरित हो जाता है; दर संभावित सेतु की
प्रकृति पर प्रबल रूप से निर्भर करती है; एक द्विनाभिकीय मध्यवर्ती
संसूचित होता है; और दर उस प्रतिस्थापन की दर से अधिक नहीं हो सकती जो
\omterm{def:b3:complex-mechanisms:labile}{विनिमयशील} साथी पर सेतु बनाने के लिए आवश्यक है।
\end{solution}

\begin{solution}{exo:b3:complex-mechanisms:9}
$(\lambda + \Delta G^\circ)^2/4\lambda$, जहाँ $4\lambda = 320$: 20.0, 11.3, 0 और \qty{5.0}{kJ/mol} (अंतिम \omterm{def:b3:complex-mechanisms:inverted}{प्रतिलोमित
क्षेत्र} में)।
\end{solution}

\begin{solution}{exo:b3:complex-mechanisms:10}
$k_{12} = \sqrt{4.0 \times \num{2.0e5} \times \num{1.0e4}} = \qty{8.9e4}{L.mol^{-1}.s^{-1}}$।
\end{solution}

\begin{solution}{exo:b3:complex-mechanisms:11}
उच्च-चक्रण $d^5$: 0; $d^7$: $1.33\,Dq$; $d^8$: $2.0\,Dq$। अपेक्षित जल-विनिमय दरें: \ce{Mn^{2+}} $>$ \ce{Co^{2+}} $>$
\ce{Ni^{2+}}, जो प्रेक्षित क्रम है।
\end{solution}

\begin{solution}{exo:b3:complex-mechanisms:12}
अत्यधिक ऊर्जाक्षेपी अभिक्रियाओं के लिए आंतरिक दर उस दर से अधिक होती है जिस पर आयन
मिलते हैं: प्रेक्षित स्थिरांक विसरण सीमा पर बना रहता है और ह्रास
छिप जाता है। एक ही अणु में दाता और ग्राही को नियत दूरी पर स्थिर करने से
कोई विसरण चरण नहीं रहता, और बढ़ते प्रेरक बल वाले ग्राहियों की एक श्रेणी ने
दर को बढ़ते, एक अधिकतम से गुज़रते और गिरते हुए दिखाया।
\end{solution}

\begin{solution}{pb:b3:complex-mechanisms:1}
\textbf{1.} $5d^8$।
\textbf{2.} $5d$ आयन का प्रबल क्षेत्र चार लिगैंडों की ओर इंगित करने वाले $d_{x^2-y^2}$ कक्षक को
बहुत ऊँचा कर देता है: आठ इलेक्ट्रॉन चार निचले कक्षकों में युग्मित हो जाते हैं।
\textbf{3.} संकुल में 16 संयोजकता इलेक्ट्रॉन और एक खुला अक्षीय स्थल है: प्रवेशी
लिगैंड पहले बँधता है, एक पंच-उपसहसंयोजी संक्रमण अवस्था से होकर।
\textbf{4.} $k_{\mathrm{obs}} = k_1 + k_2[\mathrm Y]$: एक विलायक पथ (विलायक आक्रमण, फिर Y द्वारा तीव्र प्रतिस्थापन)
और एक प्रत्यक्ष पथ।
\textbf{5.} प्रतिस्थापन में घंटों लगते हैं (भाग III): मिनटों के पैमाने पर अक्रिय, एक दिन के पैमाने पर
नहीं।
\textbf{6.} सिसप्लैटिन, क्योंकि दूसरा \ce{NH3} क्लोराइड के ट्रांस स्थित एक क्लोराइड को प्रतिस्थापित करता है
(\omterm{def:b3:complex-mechanisms:trans-effect}{ट्रांस प्रभाव} $\ce{Cl-} > \ce{NH3}$); पार्श्व उत्पाद भी बनते हैं।
\textbf{7.} आयोडाइड का \omterm{def:b3:complex-mechanisms:trans-effect}{ट्रांस प्रभाव} क्लोराइड से अधिक है, जिससे दिशा-निर्धारक
प्रभाव अधिक स्पष्ट और तीव्र हो जाता है।
\textbf{8.} $\ce{[PtI3(NH3)]-}$ में आयोडाइड के ट्रांस स्थित आयोडाइड \omterm{def:b3:complex-mechanisms:labile}{विनिमयशील} बनता है, न कि
\ce{NH3} के ट्रांस स्थित: दूसरा अमोनिया पहले के सिस स्थान पर प्रवेश करता है।
\textbf{9.} सिल्वर आयोडाइड अवक्षेपित होता है और $cis$-$\ce{[Pt(H2O)2(NH3)2]^{2+}}$ विलयन में रहता है।
\textbf{10.} क्लोराइड दोनों जल अणुओं को उनके अपने स्थानों पर प्रतिस्थापित करता है;
ऐम्मीनों को छुआ नहीं जाता।
\textbf{11.} $\ce{[Pt(NH3)4]^{2+}}$ और दो क्लोराइडों से।
\textbf{12.} इसके दोनों क्रियाशील स्थान एक-दूसरे के सम्मुख हैं: यह एक ही DNA रज्जुक के दो
पड़ोसी क्षारकों से नहीं बंध सकता, वह क्षति जिसके माध्यम से सिसप्लैटिन
कार्य करता है।
\textbf{13.} $\ln2/\num{9.0e-5} = \qty{7.7e3}{s} = \qty{2.1}{h}$।
\textbf{14.} $1 - \eu^{-\num{9.0e-5} \times 7200} = 0.48$।
\textbf{15.} जल क्लोराइड की तुलना में कहीं बेहतर निर्गामी समूह है: ऐक्वा संकुल
किसी DNA क्षारक के नाइट्रोजन द्वारा शीघ्र प्रतिस्थापित हो जाता है।
\textbf{16.} सहयोजी, जैसा सभी प्लैटिनम(II) प्रतिस्थापनों में: ऋणात्मक \omterm{def:b3:complex-mechanisms:activation-volume}{सक्रियण
आयतन}।
\textbf{17.} उच्च क्लोराइड सांद्रता साम्य को वापस
डाइक्लोरो रूप की ओर धकेलती है, जिससे औषध कोशिकाओं तक पहुँचने तक अक्षुण्ण रहती है।
\textbf{18.} $f = K/(K + [\ce{Cl-}])$।
\textbf{19.} $\num{3.0e-3}/0.103 = 0.029$।
\textbf{20.} $\num{3.0e-3}/0.0070 = 0.43$।
\textbf{21.} 15।
\textbf{22.} केवल कोशिकाओं के भीतर ही उल्लेखनीय अंश क्रियाशील ऐक्वा
संकुल के रूप में होता है।
\textbf{23.} सल्फ़र लिगैंड, जैसे ग्लूटाथायोन और प्रोटीनों के सिस्टीन और मेथियोनीन,
जो मृदु प्लैटिनम से प्रबलता से बँधते हैं।
\textbf{24.} अमोनिया प्लैटिनम से प्रबलता से बँधता है और क्लोराइड के ट्रांस स्थित है, जो
कम \omterm{def:b3:complex-mechanisms:trans-effect}{ट्रांस प्रभाव} वाला लिगैंड है: ऐम्मीन आबंध \omterm{def:b3:complex-mechanisms:labile}{विनिमयशील} नहीं बनते।
\textbf{25.} \textbf{कोशिका के भीतर औषध प्लाज़्मा की तुलना में लगभग 15 गुना अधिक
क्रियाशील ऐक्वा रूप में होती है।}
\end{solution}
